% !TEX root = ../main.tex
%% Corrigé du devoir surveillé n° 4 — vérifié par verifications/eval_ds.py.

\corrigeSujetEval{DS4}{Devoir surveillé \numero4}

\begin{corrigeEval}{6 points}{Développements limités}
\begin{questions}
  \item On compose avec $u=\sin x=x-\frac{x^{3}}6+o(x^{4})$, qui tend vers
    $0$ : $u^{2}=x^{2}-\frac{x^{4}}3+o(x^{4})$, $u^{3}=x^{3}+o(x^{4})$,
    $u^{4}=x^{4}+o(x^{4})$. Dans
    $\mathrm{e}^{u}=1+u+\frac{u^{2}}2+\frac{u^{3}}6+\frac{u^{4}}{24}+o(u^{4})$,
    les termes en $x^{3}$ se compensent et ceux en $x^{4}$ donnent
    $-\frac16+\frac1{24}=-\frac18$ :
    \[\res{\mathrm{e}^{\sin x}=1+x+\frac{x^{2}}{2}-\frac{x^{4}}{8}+o(x^{4})}.\]
    \pts{2 pts}
  \item $\mathrm{e}^{x}=1+x+\frac{x^{2}}2+\frac{x^{3}}6+o(x^{3})$, donc
    $\mathrm{e}^{\sin x}-\mathrm{e}^{x}=-\frac{x^{3}}6+o(x^{3})$ : la limite
    vaut $\res{-\dfrac16}$. \pts{1 pt}
  \item Pour $x>0$, avec $\sqrt{1+h}=1+\frac h2-\frac{h^{2}}8+o(h^{2})$ et
    $h=\frac2x\to0$ :
    \[\sqrt{x^{2}+2x}=x\sqrt{1+\frac2x}
      =x\Bigl(1+\frac1x-\frac{1}{2x^{2}}+o\Bigl(\frac{1}{x^{2}}\Bigr)\Bigr)
      =x+1-\frac{1}{2x}+o\Bigl(\frac1x\Bigr).\]
    Asymptote $\res{y=x+1}$ ; la courbe est \res{\text{au-dessous}}, car
    l'écart est équivalent à $-\frac{1}{2x}<0$. \pts{1,5 pt}
  \item On réduit au même dénominateur :
    \[\frac{1}{\sin^{2}x}-\frac1{x^{2}}=\frac{x^{2}-\sin^{2}x}{x^{2}\sin^{2}x}
      =\frac{\frac{x^{4}}{3}+o(x^{4})}{x^{4}+o(x^{4})}\longrightarrow\res{\frac13}.\]
    \pts{1,5 pt}
\end{questions}
\begin{center}\begin{minipage}{\linewidth}\centering
\begin{tikzpicture}[scale=0.6,line cap=round]
  \draw[otMuted,line width=0.5pt,-{Stealth[length=1.5mm]}] (-0.3,0) -- (6.4,0);
  \draw[otMuted,line width=0.5pt,-{Stealth[length=1.5mm]}] (0,-0.3) -- (0,7.4);
  \draw[otGreenDark,line width=0.9pt,dash pattern=on 3pt off 1.6pt] (-0.3,0.7) -- (6.2,7.2);
  \draw[otNavy,line width=1.2pt] plot[domain=0:6.2,samples=90] ({\x},{sqrt(\x*\x+2*\x)});
  \node[font=\scriptsize,text=otGreenDark,anchor=south east] at (4.5,5.5) {$y=x+1$};
  \node[font=\scriptsize,text=otNavy,anchor=north west] at (4.6,5.2) {$y=\sqrt{x^{2}+2x}$};
\end{tikzpicture}\\*[0.5mm]
\legende{L'asymptote $y=x+1$ en $+\infty$ : la courbe reste au-dessous,
l'écart étant de l'ordre de $\frac{1}{2x}$.}
\end{minipage}\end{center}
\end{corrigeEval}

\begin{corrigeEval}{4 points}{Suite implicite}
\begin{questions}
  \item $f_n$ est continue et strictement croissante sur $\intFF01$, avec
    $f_n(0)=-1<0<1=f_n(1)$ : par le théorème de la bijection, elle s'annule
    en un unique $u_n\in\intOO01$. Pour $n=2$, $u_2^{2}+u_2-1=0$ et
    $u_2>0$ : $\res{u_2=\dfrac{\sqrt5-1}{2}}$. \pts{1 pt}
  \item Comme $u_n^{n}=1-u_n$ :
    \[f_{n+1}(u_n)=u_n^{n+1}+u_n-1=u_n^{n+1}-u_n^{n}=u_n^{n}(u_n-1)<0=f_{n+1}(u_{n+1}),\]
    et $f_{n+1}$ est croissante : $u_n<u_{n+1}$. Croissante et majorée par
    $1$, $(u_n)$ converge vers $\ell\leqslant1$. Si $\ell<1$, alors
    $0<1-u_n=u_n^{n}\leqslant\ell^{n}\to0$, d'où $\ell=1$ : contradiction.
    Donc $\res{u_n\to1}$. \pts{1,5 pt}
  \item $u_n^{n}=1-u_n=v_n$ ; en prenant le logarithme,
    $n\ln(1-v_n)=\ln v_n$. Comme $v_n\to0$, $\ln(1-v_n)\sim-v_n$, d'où
    $n\,v_n\sim-\ln v_n\to+\infty$. En prenant le logarithme de cet
    équivalent :
    \[\ln n+\ln v_n=\ln(-\ln v_n)+o(1)=o(\ln v_n),\]
    donc $\ln v_n\sim-\ln n$, puis $n\,v_n\sim\ln n$ :
    $\res{v_n\sim\dfrac{\ln n}{n}}$. \pts{1,5 pt}
\end{questions}
\begin{center}\begin{minipage}{\linewidth}\centering
\begin{tikzpicture}[x=6cm,y=2.2cm,line cap=round]
  \draw[otMuted,line width=0.5pt,-{Stealth[length=1.5mm]}] (0,0) -- (1.08,0);
  \draw[otMuted,line width=0.5pt,-{Stealth[length=1.5mm]}] (0,-1.08) -- (0,1.15);
  \foreach \n/\c in {2/40,4/60,8/80,16/100}
    \draw[otNavy!\c,line width=0.9pt] plot[domain=0:1,samples=120] ({\x},{\x^\n+\x-1});
  \foreach \r/\l in {0.618/2,0.7245/4,0.8117/8,0.8772/16}{
    \fill[otGreenDark] (\r,0) circle (1.5pt);
    \node[font=\scriptsize,text=otGreenDark,anchor=north west,inner sep=1pt] at (\r,-0.03) {$u_{\l}$};}
  \node[font=\scriptsize,anchor=north] at (1,-0.03) {$1$};
  \node[font=\scriptsize,anchor=east] at (0,1) {$1$};
  \node[font=\scriptsize,anchor=east] at (0,-1) {$-1$};
  \node[font=\scriptsize,text=otNavy!40,anchor=south east] at (0.5,-0.25) {$n=2$};
  \node[font=\scriptsize,text=otNavy,anchor=north west] at (0.36,-0.68) {$n=16$};
\end{tikzpicture}\\*[0.5mm]
\legende{Les courbes de $f_n$ pour $n=2,4,8,16$ : leurs racines $u_n$ (en vert)
croissent vers $1$, avec $1-u_n\sim\frac{\ln n}{n}$.}
\end{minipage}\end{center}
\end{corrigeEval}

\begin{corrigeEval}{5 points}{Espaces vectoriels}
\begin{questions}
  \item $\varphi$ est linéaire (évaluation et dérivation le sont) ;
    $\varphi(1)=(1,0)$ et $\varphi(X-1)=(0,1)$ : l'image contient une base
    de $\R^{2}$, $\varphi$ est surjective. $F=\operatorname{Ker}\varphi$ est
    un sous-espace vectoriel, et le théorème du rang donne
    $\dim F=4-2=\res{2}$. \pts{2 pts}
  \item $P(1)=P'(1)=0$ équivaut à « $1$ est racine au moins double de
    $P$ », soit $(X-1)^{2}\mid P$ : $F=(X-1)^{2}\,\R_1[X]$, de base
    $\res{\bigl((X-1)^{2},\,X(X-1)^{2}\bigr)}$ (famille libre de degrés
    distincts, de cardinal $2=\dim F$). \pts{1,5 pt}
  \item Un $P\in F\cap G$ est de degré au plus $1$ avec $1$ pour racine
    double : $P=0$. Comme $\dim F+\dim G=2+2=4$, $E=F\oplus G$. La division
    euclidienne de $X^{3}$ par $(X-1)^{2}$ (DS 3) donne le reste $3X-2$ :
    \[\res{X^{3}=\underbrace{(X-1)^{2}(X+2)}_{=\,X^{3}-3X+2\,\in\,F}+\underbrace{(3X-2)}_{\in\,G}}.\]
    \pts{1,5 pt}
\end{questions}
\end{corrigeEval}

\begin{corrigeEval}{5 points}{Projecteur et matrices}
\begin{questions}
  \item $e_1\notin\mathcal{P}$ ($1+0+0\neq0$), donc
    $\mathcal{P}\cap D=\ens0$ ; et $\dim\mathcal{P}+\dim D=2+1=3$ :
    $\res{\R^{3}=\mathcal{P}\oplus D}$. \pts{1 pt}
  \item On écrit $(x,y,z)=p(x,y,z)+t\,e_1$ avec $p(x,y,z)\in\mathcal{P}$ :
    $(x-t)+y+z=0$, soit $t=x+y+z$ et $p(x,y,z)=(-y-z,\ y,\ z)$. Donc
    \[M=\res{\begin{pmatrix}0&-1&-1\\0&1&0\\0&0&1\end{pmatrix}},\qquad M^{2}=M,\qquad
      \operatorname{tr}M=2=\operatorname{rg}M,\]
    conformément au résultat du cours : la trace d'un projecteur est égale
    à son rang. \pts{1,5 pt}
  \item Dans une base adaptée à $\R^{3}=\mathcal{P}\oplus D$ : $u_1=(1,-1,0)$,
    $u_2=(1,0,-1)$ (base de $\mathcal{P}$) et $u_3=e_1$. Avec
    \[Q=\begin{pmatrix}1&1&1\\-1&0&0\\0&-1&0\end{pmatrix},\qquad
      \res{Q^{-1}MQ=\operatorname{diag}(1,1,0)}.\]
    \pts{1,5 pt}
  \item $s=2p-\mathrm{id}$ a pour matrice
    $\res{\begin{pmatrix}-1&-2&-2\\0&1&0\\0&0&1\end{pmatrix}}$, et
    $s^{2}=4p^{2}-4p+\mathrm{id}=\mathrm{id}$ car $p^{2}=p$. \pts{1 pt}
\end{questions}
\end{corrigeEval}
